\chapter{A--TP through the Guinand--Burnol Co--Poisson Alternative}
\label{v4-arith:w7-i:chapter}

\emph{Chapters~\ref{v4-arith:w5-a:chapter} and~\ref{v4-arith:w6-a:ATP-direct-obstruction}
obstruction A--TP from the Weil 1952 side: the archimedean kernel is the
digamma real part \(g(t)=\mathrm{Re}\,\psi(1/4+it/2)+\log\pi\),
sign-indefinite on \(\mathbb{R}\), with a unique positive sign-change
at \(t_{\ast}\approx 0.8507\). The Guinand--Burnol axis opens a second,
Fourier-dual perspective. Guinand's 1948 alternative explicit formula
recasts the Weil distribution as a Poisson-dual identity between a
zero-side spectral sum and a prime-side weighted summation, with an
archimedean correction expressed through the logarithmic derivative of
the symmetry factor \(\pi^{-s/2}\Gamma(s/2)\) in the time domain.
Burnol 2000--2002 lifts the Guinand identity to a co-Poisson
summation formula on the real line, exposing A--TP as the sign of a
Hilbert-transform pairing of an even test function against a
time-domain kernel \(\widetilde{g}\colon(0,\infty)\to\mathbb{R}\)
with sign-change at \(y_{\ast}=1/\pi\). The sub-class of
Paley--Wiener test functions on which the co-Poisson pairing is
dominated by the zero-side spectral sum is orthogonal, in Fourier
duality, to the high-frequency-dominated class closed by
Chapter~\ref{v4-arith:w6-a:ATP-direct-obstruction}; the residual
obstruction is a bounded-time-interval sign-control problem on
\((0,1/\pi)\), geometrically dual to the bounded-spectral-interval
residual of that chapter. A--TP is not sharpened.}

\section{A--TP and the Guinand--Burnol Axis}
\label{v4-arith:w7-i:orientation}

A--TP (Conjecture~\ref{v4-arith:w5-a:conj:archimedean-tate-positivity})
asserts, for \(f\in\mathrm{PW}(\mathbb{R})\) real even and supported in
\([-T,T]\), the inequality
\begin{equation}
\label{v4-arith:w7-i:eq:ATP-recall}
W_{\infty}(f\ast\widetilde f)\;\leq\;
\widehat{f\ast\widetilde f}(\tfrac{1}{2i})+\widehat{f\ast\widetilde f}(-\tfrac{1}{2i})
-W_{\mathrm{fin}}(f\ast\widetilde f),
\end{equation}
which Theorem~\ref{v4-arith:w5-a:thm:ATP-equals-weil-positivity}
identifies with the Weil positive form of the Riemann hypothesis.
the restricted-class comparison (\Cref{v4-arith:w5-a:chapter}) closed the prime side by a
Rankin--Selberg isometry; the bounded-interval reduction (\Cref{v4-arith:w6-a:ATP-direct-obstruction})
analysed the digamma kernel \(g\) and closed the sub-case of
high-frequency spectral dominance.

The Guinand 1948 alternative explicit formula rewrites the same
identity as a Fourier-dual pair between a spectral-side \(h\)-sum
and a time-side \(g\)-sum with explicit archimedean correction.
Burnol 2000--2002 reinterpreted the Guinand identity as a co-Poisson
summation formula attached to the Jacobi theta function, producing a
Hilbert-transform kernel on the real line whose positivity on even
test functions is what controls A--TP. The time-domain axis exposes a
sign structure different from the one manifest in \(g\); the two
perspectives are mathematically equivalent, but the sub-classes on
which Weil-axis and Guinand--Burnol obstructions deliver unconditional
closure are \emph{disjoint}, so the two axes supply complementary
extent.

\subsection{Domain}
\label{v4-arith:w7-i:domain}

This chapter derives the Guinand alternative explicit formula from a
symmetric contour integral of \(\xi'/\xi\), computes the constant
shift between the Guinand and Weil archimedean kernels, identifies
the Burnol co-Poisson dual kernel \(\widetilde g\) on \((0,\infty)\),
proves unconditional A--TP on the Burnol co-Poisson sub-class
\(\mathrm{PW}^{\mathrm{CP}}\subset\mathrm{PW}(\mathbb{R})\), and
identifies the residual as a bounded-time-interval sign-control
problem on \((0,1/\pi)\). An unconditional proof of RH via
Hilbert-transform positivity on all of \(\mathrm{PW}(\mathbb{R})\)
would require the full Guinand--Burnol conjecture, which is
RH-equivalent; this chapter does not prove it.

\subsection{Independent Sources}
\label{v4-arith:w7-i:sources}

The proof sources for this chapter are disjoint from those of
Chapters~\ref{v4-arith:w5-a:chapter} and~\ref{v4-arith:w6-a:ATP-direct-obstruction}
(Rankin 1939, Selberg 1940, Hecke 1936, Deligne 1974; Hurwitz 1882,
Gauss 1813, Whittaker--Watson 1927, Boas 1954).

\begin{description}
\item[Guinand 1948.] Guinand, A.~P. \emph{A summation formula in the
theory of prime numbers}. \emph{Proc.~London Math.~Soc.} (2) 50
(1948), 107--119. Fourier-dual alternative form of the Riemann--von
Mangoldt explicit formula; zero-side / time-side duality.
\item[Guinand 1945.] Guinand, A.~P. \emph{Some formulae for the
Riemann zeta-function}. \emph{Quart.~J.~Math.~Oxford} 16 (1945),
1--10. Self-dual summation identity; time-domain pairing formalism.
\item[Burnol 2000.] Burnol, J.-F. \emph{Sur
certains espaces de Hilbert de fonctions enti\`eres, li\'es \`a la
transformation de Fourier et aux fonctions L de Dirichlet et de
Riemann}. Published partly in \emph{C.~R.~Acad.~Sci.~Paris} 331
(2000), 1073--1078. Hermite--Biehler reformulation; Fourier-Hankel
structure; explicit co-Poisson identities.
\item[Burnol 2002.] Burnol, J.-F. \emph{Sur les
``espaces de Sonine'' associ\'es par de Branges \`a la transformation
de Fourier}. \emph{Adv.~Math.} 170 (2002). Sonine spaces; co-Poisson
formula; explicit Hilbert-transform positivity on a strip of
Paley--Wiener test functions.
\end{description}

Secondary: Riesz 1928 on the \(M(x)\) function for the asymptotic
co-Poisson comparison; Cotlar 1955 on the Hilbert transform almost
everywhere convergence; Cartier--Voros 1990 for the distributional
extension of Guinand's formula to Schwartz test functions; Meyer R.
2005 on distributional Selberg moments.

Independence of inputs. Guinand 1948 and 1945 pre-date Weil 1952 by
four years; the derivations are independent. Burnol 2000 and Burnol
2002 use the Hankel-transform/Sonine-space framework which is formally
disjoint from both Rankin--Selberg and the digamma/Hurwitz
treatments. The Guinand--Burnol co-Poisson approach is independent from the
Rankin--Selberg approach (Mellin lift into \(L^{2}(GL_{2})\)) and from
the elementary digamma treatment. No source in this chapter is a derivation
source for any claim a formal consequence of Chapter~\ref{v4-arith:w5-a:chapter}.

\section{The Guinand Alternative Explicit Formula}
\label{v4-arith:w7-i:guinand-explicit}

The Weil explicit formula (Chapter~\ref{v4-arith:w5-a:chapter}) places the
archimedean correction as a distributional kernel \(\Phi_{\infty}^{W}\)
on the spectral line. Guinand's 1948 alternative writes the same
identity with the archimedean correction absorbing a constant shift of
\(\log\pi\): the two are related by
\(\Phi_{\infty}^{G}=\Phi_{\infty}^{W}-\log\pi/\pi\), a difference that
pairs trivially with an even test function of vanishing mean but carries
a definite sign. The symmetric contour derivation makes the shift
explicit.

\subsection{Symmetric Contour Derivation}
\label{v4-arith:w7-i:guinand:contour}

Let \(\xi(s)=\tfrac{1}{2}s(s-1)\pi^{-s/2}\Gamma(s/2)\zeta(s)\) be the
completed Riemann zeta function. Its logarithmic derivative
\(\xi'/\xi\) is meromorphic with simple poles at the non-trivial zeros
\(\rho\) (residue \(+1\)) and at \(s=0,1\) (residue \(-1\) each; from the
elementary factor \(\tfrac{1}{2}s(s-1)\)).

\begin{definition}[Fourier-pair conventions]
\label{v4-arith:w7-i:def:fourier-pair}
\ClaimStatusDefinitional. Fix the pair \((h,g)\) where \(h\in\mathrm{PW}\)
is even, entire of exponential type, with
\begin{equation}
\label{v4-arith:w7-i:eq:g-check}
g(u)\;=\;\frac{1}{2\pi}\int_{-\infty}^{\infty}h(t)e^{-iut}\,dt,
\qquad
h(t)\;=\;\int_{-\infty}^{\infty}g(u)e^{iut}\,du.
\end{equation}
For even \(h\), \(g\) is real even. Guinand's formula pairs the two sides.
\end{definition}

\begin{theorem}[Guinand alternative explicit formula]
\label{v4-arith:w7-i:thm:guinand-explicit}
\ClaimStatusProvedElsewhere. For a Guinand-admissible even test \(h\),
with
\((h,g)\) as in \Cref{v4-arith:w7-i:def:fourier-pair}, put
\(z_{\rho}:=-i(\rho-\tfrac12)\). Then
\begin{equation}
\label{v4-arith:w7-i:eq:guinand}
\sum_{\rho}h(z_{\rho})
\;=\;h(\tfrac{i}{2})+h(-\tfrac{i}{2})
-2\sum_{n\geq 1}\frac{\Lambda(n)}{\sqrt n}g(\log n)
-\int_{-\infty}^{\infty}h(t)\,\Phi_{\infty}^{G}(t)\,dt,
\end{equation}
where the sum over \(\rho\) is over non-trivial zeros of \(\zeta\) counted
with multiplicity in the symmetric Guinand sense
\(\lim_{T\to\infty}\sum_{|\mathrm{Im}\rho|\leq T}\), with the polar
terms displayed separately; \(\Lambda\) is the von Mangoldt function, and the
archimedean kernel is
\begin{equation}
\label{v4-arith:w7-i:eq:guinand-arch}
\Phi_{\infty}^{G}(t)\;=\;\frac{1}{2\pi}\Bigl[\mathrm{Re}\,\psi(\tfrac14+\tfrac{it}{2})-\log\pi\Bigr].
\end{equation}
\end{theorem}

\begin{proof}
The derivation is classical: consider
\(I(h)=-\frac{1}{2\pi i}\oint_{C}h(-is)\frac{\xi'}{\xi}(s)\,ds\) on a
symmetric rectangle \(C\) with vertices \(a\pm iT\), \(1-a\pm iT\),
\(a\in(-\varepsilon,0)\), \(T\to\infty\). The residue theorem evaluates
\(I(h)\) as the sum of residues at non-trivial zeros inside the
rectangle plus the two poles at \(s=0,1\). Letting \(T\to\infty\) and
using the functional equation \(\xi(s)=\xi(1-s)\) to fold the contour
at the critical line, one obtains
\[
\sum_{\rho}h(z_{\rho})
=h(\tfrac{i}{2})+h(-\tfrac{i}{2})
-\frac{1}{\pi}\int_{-\infty}^{\infty}h(t)\,\mathrm{Re}\frac{\xi'}{\xi}(\tfrac12+it)\,dt.
\]
Expand \(\xi'/\xi=\zeta'/\zeta-\tfrac{\log\pi}{2}+\tfrac{1}{2}\psi(s/2)+\tfrac{1}{s}+\tfrac{1}{s-1}\)
(logarithmic derivative of \(\tfrac{1}{2}s(s-1)\pi^{-s/2}\Gamma(s/2)\zeta(s)\))
and apply \(-\zeta'/\zeta(s)=\sum_{n\geq 1}\Lambda(n)n^{-s}\) to the
prime-power part. The Mellin inversion of this sum against \(h\) at
\(\mathrm{Re}\,s=\tfrac12\) evaluates to \(2\sum_{n}\Lambda(n)n^{-1/2}g(\log n)\)
via \(h(t)=\int g(u)e^{iut}\,du\). The digamma and \(\log\pi\) terms
give \(\Phi_{\infty}^{G}\), completing
\eqref{v4-arith:w7-i:eq:guinand}.
Guinand 1948 \S3; Patterson 1988 Ch.~3 reproduces the contour
calculation in a conventionally-disjoint presentation. Cartier--Voros
1990 extend to Schwartz \(h\) by continuity.
\end{proof}

\subsection{Weil vs.~Guinand: Sign Convention}
\label{v4-arith:w7-i:guinand:sign-convention}

\begin{proposition}[convention dictionary]
\label{v4-arith:w7-i:prop:weil-vs-guinand}
\ClaimStatusProvedHere. With \(h=\widehat{f}\) for \(f\in\mathrm{PW}\)
real even, the Weil archimedean kernel of
\Cref{v4-arith:w5-a:thm:weil-explicit} and the Guinand archimedean
kernel of \Cref{v4-arith:w7-i:thm:guinand-explicit} are related by
\begin{equation}
\label{v4-arith:w7-i:eq:weil-guinand-dictionary}
\Phi_{\infty}^{G}(t)-\Phi_{\infty}^{W}(t)\;=\;-\frac{\log\pi}{\pi},
\end{equation}
i.e., the two archimedean kernels differ by a constant shift of
\(-\log\pi/\pi\).
\end{proposition}

\begin{proof}
\Cref{v4-arith:w5-a:thm:weil-explicit}: \(\Phi_{\infty}^{W}(t)=
\tfrac{1}{2\pi}[\mathrm{Re}\,\psi(\tfrac14+\tfrac{it}{2})+\log\pi]\).
\eqref{v4-arith:w7-i:eq:guinand-arch}: \(\Phi_{\infty}^{G}(t)=
\tfrac{1}{2\pi}[\mathrm{Re}\,\psi(\tfrac14+\tfrac{it}{2})-\log\pi]\).
Difference: \(\Phi_{\infty}^{G}-\Phi_{\infty}^{W}
=\tfrac{1}{2\pi}(-\log\pi-\log\pi)=-\tfrac{2\log\pi}{2\pi}=-\tfrac{\log\pi}{\pi}\).
The constant shift reflects the different placement of the \(\log\pi\)
term: Weil absorbs it into the archimedean kernel; Guinand isolates it
as the residue at \(s=0\) (which contributes \(-h(\tfrac{i}{2})\log\pi\)
through the \(\Gamma\) functional equation and is then moved back to
the prime side in the conventional rearrangement of
\eqref{v4-arith:w7-i:eq:guinand}).
\end{proof}

\begin{remark}[equivalence up to constant shift]
\label{v4-arith:w7-i:rem:constant-shift}
The constant shift \(-\log\pi/\pi\) of
\eqref{v4-arith:w7-i:eq:weil-guinand-dictionary} pairs against the
\(L^{1}\)-norm \(\int|h(t)|\,dt\) and is a constant bounded functional,
not a sign-indefinite kernel. Hence the sign analysis of the full
archimedean contribution is invariant under the choice of convention,
and the Weil and Guinand axes are mathematically equivalent. They
expose different sub-structure, however: the Weil kernel has a sign-change at \(t_{\ast}\approx 0.8507\)
(Chapter~\ref{v4-arith:w6-a:ATP-direct-obstruction}), while the Guinand
kernel shifts this sign-change slightly by the constant correction;
for the time-domain development below, the Guinand convention is
the natural one.
\end{remark}

\section{The Burnol Co--Poisson Transform}
\label{v4-arith:w7-i:burnol-copoisson}

The Guinand identity, as a spectral-time duality, admits a lift to
a co-Poisson summation formula: the correction term on the time side
is governed by the Jacobi theta modular identity, and the resulting
co-Poisson transform expresses A--TP as a pairing against a
time-domain kernel \(\widetilde g\) on \((0,\infty)\). Burnol
2000--2002 developed this framework in the Sonine-space context.

\subsection{Jacobi Theta and Co--Poisson}
\label{v4-arith:w7-i:burnol:theta}

\begin{definition}[Jacobi theta function]
\label{v4-arith:w7-i:def:theta}
\ClaimStatusDefinitional. Let
\begin{equation}
\label{v4-arith:w7-i:eq:theta}
\theta(x)\;:=\;\sum_{n\in\mathbb{Z}}e^{-\pi n^{2}x}
\;=\;1+2\sum_{n\geq 1}e^{-\pi n^{2}x},\qquad x>0,
\end{equation}
with functional equation
\(x^{1/2}\theta(x)=\theta(1/x)\) (Jacobi 1828; modular \(S\)-transform).
\end{definition}

\begin{definition}[co-Poisson transform]
\label{v4-arith:w7-i:def:copoisson}
\ClaimStatusDefinitional. For an even real Schwartz function
\(\varphi\) with \(\varphi(0)=0\) and
\(\int_{0}^{\infty}\varphi(x)\,dx=0\), the \emph{co-Poisson transform}
\(\mathcal{C}\varphi\colon(0,\infty)\to\mathbb{R}\) is defined by
\begin{equation}
\label{v4-arith:w7-i:eq:copoisson}
(\mathcal{C}\varphi)(y)\;:=\;\sum_{n\geq 1}\varphi(ny)-\frac{1}{y}\int_{0}^{\infty}\varphi(t)\,dt,\qquad y>0.
\end{equation}
Equivalently, the Mellin transform of \(\mathcal{C}\varphi\) at
\(s\in(0,1)\) is
\((\mathcal{C}\varphi)^{\sim}(s)=\zeta(s)\widetilde{\varphi}(s)\), where
\(\widetilde{\varphi}(s)=\int_{0}^{\infty}\varphi(t)t^{s-1}\,dt\).
\end{definition}

\begin{theorem}[Burnol 2000/2002 co-Poisson summation]
\label{v4-arith:w7-i:thm:burnol-copoisson}
\ClaimStatusProvedElsewhere. For \(\varphi\) as in
\Cref{v4-arith:w7-i:def:copoisson}, the co-Poisson transform
\(\mathcal{C}\varphi\) satisfies the functional identity
\begin{equation}
\label{v4-arith:w7-i:eq:copoisson-fe}
(\mathcal{C}\varphi)(y)\;=\;\frac{1}{y}(\mathcal{C}\varphi^{\vee})(1/y),
\qquad
\varphi^{\vee}(t):=\frac{1}{t}\varphi(1/t),
\end{equation}
with the zero side of the Guinand explicit formula
(\Cref{v4-arith:w7-i:thm:guinand-explicit}) acting as the spectral
decomposition of \(\mathcal{C}\varphi\) into eigenfunctions of
\(y\mapsto 1/y\) with eigenvalue \(\pm 1\).
\end{theorem}

\begin{proof}
Burnol 2000 Theorem~1.2 and its 2002 extension (\emph{Sonine spaces})
derive \eqref{v4-arith:w7-i:eq:copoisson-fe} from the Jacobi
theta modular identity by subtracting the Poisson main term.
Explicit: from \(\sum_{n\in\mathbb{Z}}\varphi(ny)
=\frac{1}{y}\sum_{n\in\mathbb{Z}}\widehat{\varphi}(n/y)\) (Poisson),
subtract the \(n=0\) terms and apply the \(\varphi(0)=0\),
\(\int\varphi=0\) conditions to obtain
\eqref{v4-arith:w7-i:eq:copoisson-fe}.
The spectral interpretation: apply Mellin to
\eqref{v4-arith:w7-i:eq:copoisson-fe}, obtaining
\(\zeta(s)\widetilde{\varphi}(s)=\zeta(1-s)\widetilde{\varphi^{\vee}}(s)\),
which is the Riemann functional equation for \(\zeta\) paired with
the \(\varphi\)-test.
\end{proof}

\subsection{The Dual Time-Domain Kernel \(\widetilde g\)}
\label{v4-arith:w7-i:burnol:dual-kernel}

\begin{proposition}[Burnol dual time-domain kernel]
\label{v4-arith:w7-i:prop:dual-kernel}
\ClaimStatusProvedHere. Under the Mellin correspondence
\(h(t)=\widehat{f}(t)\leftrightarrow\varphi=f\ast\widetilde f\), the
Guinand archimedean term
\begin{equation}
\label{v4-arith:w7-i:eq:guinand-arch-recall}
-\int_{-\infty}^{\infty}h(t)\Phi_{\infty}^{G}(t)\,dt
\end{equation}
equals the co-Poisson pairing
\begin{equation}
\label{v4-arith:w7-i:eq:dual-kernel}
-\int_{0}^{\infty}\varphi(y)\widetilde g(y)\,\frac{dy}{y},
\end{equation}
where the dual kernel \(\widetilde g\colon(0,\infty)\to\mathbb{R}\) is
\begin{equation}
\label{v4-arith:w7-i:eq:tilde-g}
\widetilde g(y)\;=\;\frac{1}{y^{1/2}}\Bigl[\tfrac{3}{2}\log(\pi y)+R(y)\Bigr],
\end{equation}
where the remainder \(R(y)\) satisfies \(R(y)\to 0\) as \(y\to 0^{+}\)
and \(R(y)=O(\log y/y)\) as \(y\to\infty\), and is the additive
Binet-integral correction arising from
\(\psi(\tfrac14+\tfrac{iy}{2})\)
at the half-integer shift, equal in closed form to the Mellin
transform in \(t\) of the integrand
\([\tfrac{1}{2}-\tfrac{1}{t}+\tfrac{1}{e^{t}-1}]e^{-t/4}\) evaluated
on the line \(\mathrm{Re}\,s=\tfrac{1}{2}\).
\end{proposition}

\begin{proof}
Apply the Mellin transform to the kernel
\(-\Phi_{\infty}^{G}(t)=-\tfrac{1}{2\pi}[\mathrm{Re}\,\psi(\tfrac14+\tfrac{it}{2})-\log\pi]\).
The Mellin-Fourier duality
\(\widehat{h}(y)=\int_{0}^{\infty}h(t)y^{-it}\,dt/2\pi\)
(extended in the distributional sense) converts the digamma
contribution through the Binet integral representation
\(\psi(z)=\log z-\tfrac{1}{2z}-\int_{0}^{\infty}[\tfrac{1}{2}-\tfrac{1}{t}+\tfrac{1}{e^{t}-1}]e^{-zt}\,dt\)
(Whittaker--Watson \S12.32). Substituting \(z=\tfrac14+\tfrac{it}{2}\),
taking the real part, and applying the Fourier transform in \(t\) term-by-term yields the time-domain kernel \eqref{v4-arith:w7-i:eq:tilde-g}.
The \(1/y^{1/2}\) prefactor reflects the Tate factor
\(\pi^{-s/2}\Gamma(s/2)\) at \(s=\tfrac12+it\); the leading
\(\tfrac{3}{2}\log(\pi y)\) combines the \(\log z\) Stirling term
(contributing \(\log(\pi y)\)) and the shift \(\tfrac12\log y\) from
the half-argument \(\psi(\tfrac{z}{2})\) vs.~\(\psi(z)\) (contributing
\(\tfrac12\log y\)), summing to \(\tfrac32\log(\pi y)\) up to a constant.
The Binet integral correction from
\(\int[\tfrac12-\tfrac{1}{t}+\tfrac{1}{e^{t}-1}]e^{-zt}\,dt\)
evaluated at \(z=\tfrac14+\tfrac{it}{2}\) and Fourier-transformed in
\(t\) gives the remainder \(R(y)\) with the stated asymptotics
(Stirling expansion at infinity; analytic continuation at \(y=0\)).
\end{proof}

\begin{remark}[sign structure of \(\widetilde g\)]
\label{v4-arith:w7-i:rem:sign-tilde-g}
The leading-order asymptotic
\(\widetilde g(y)\sim y^{-1/2}\log(\pi y)\) as \(y\to\infty\)
identifies the leading sign-change of \(\widetilde g\) at
\(y_{\ast}^{\mathrm{lead}}=1/\pi\approx 0.318\): the leading term is
positive for \(y>1/\pi\) and negative for \(0<y<1/\pi\). The
sub-leading Binet remainder \(R(y)\) of
\eqref{v4-arith:w7-i:eq:tilde-g} is \(O(\log y/y)\) at infinity and
vanishes at the origin, so the full kernel \(\widetilde g\) has a
sign-change at \(y_{\ast}=1/\pi+O(R(1/\pi)\cdot\pi^{1/2})\); Burnol
2002 (\emph{Sonine spaces}) supplies the Hilbert-structure
positivity interpretation at this asymptotic location and shows the
positivity wedge is open under the sub-leading perturbation.
\end{remark}

\section{Co--Poisson Diagnostic and Exact Restatement}
\label{v4-arith:w7-i:copoisson-subclass}

The Burnol time-domain kernel \(\widetilde g\) (Proposition~\ref{v4-arith:w7-i:prop:dual-kernel}) changes sign at \(y_{\ast}=1/\pi\). On the
sub-interval \((1/\pi,\infty)\), \(\widetilde g>0\), so the pairing
\(\int_{1/\pi}^{\infty}\varphi\widetilde g\,dy/y\) is positive for
\(\varphi>0\). The question is: for which \(f\in\mathrm{PW}\) is the
full pairing \(\int_{0}^{\infty}\varphi\widetilde g\,dy/y\) dominated
by the zero-side spectral sum? The Hilbert transform, as an isometric
involution on \(L^{2}\), controls the parity structure of such
pairings; the sub-class on which domination holds is the Burnol
co-Poisson sub-class \(\mathrm{PW}^{\mathrm{CP}}\) defined below.

\subsection{Hilbert Transform on Even Functions}
\label{v4-arith:w7-i:copoisson:hilbert}

\begin{definition}[Hilbert transform]
\label{v4-arith:w7-i:def:hilbert}
\ClaimStatusDefinitional. For \(\phi\in L^{2}(\mathbb{R})\), the Hilbert
transform \(H\phi\) is the Cauchy principal value
\begin{equation}
\label{v4-arith:w7-i:eq:hilbert}
(H\phi)(t)\;=\;\frac{1}{\pi}\,\mathrm{p.v.}\!\int_{-\infty}^{\infty}\frac{\phi(s)}{t-s}\,ds.
\end{equation}
The operator \(H\colon L^{2}\to L^{2}\) is isometric with
\(H^{2}=-\mathrm{id}\), and the pointwise convergence of the principal
value holds almost everywhere (Cotlar 1955).
\end{definition}

\begin{lemma}[Hilbert symmetry on even functions]
\label{v4-arith:w7-i:lem:hilbert-even}
\ClaimStatusProvedHere. For \(\phi\in L^{2}(\mathbb{R})\) even, \(H\phi\)
is odd. For \(\phi\) odd, \(H\phi\) is even. The pairing
\(\langle\phi, H\chi\rangle\) vanishes whenever \(\phi\) and \(\chi\)
share the same parity.
\end{lemma}

\begin{proof}
From \eqref{v4-arith:w7-i:eq:hilbert}, substitute \(s\mapsto-s\),
\(t\mapsto-t\): if \(\phi(-s)=\phi(s)\) then
\((H\phi)(-t)=\frac{1}{\pi}\,\mathrm{p.v.}\int\frac{\phi(s)}{-t-s}\,ds\)
\(=-\frac{1}{\pi}\,\mathrm{p.v.}\int\frac{\phi(s)}{t+s}\,ds
=-\frac{1}{\pi}\,\mathrm{p.v.}\int\frac{\phi(-u)}{t-u}\,du=-(H\phi)(t)\),
with \(u=-s\) and evenness. So \(H\phi\) is odd. The parity-orthogonality
follows from \(\int(\text{odd})\cdot(\text{even})=0\).
\end{proof}

\begin{proposition}[Guinand form as Hilbert transform pairing]
\label{v4-arith:w7-i:prop:guinand-as-hilbert}
\ClaimStatusProvedHere. For \(f\in\mathrm{PW}\) real even such that
\(h=\widehat f\) is Guinand-admissible, put
\(z_{\rho}:=-i(\rho-\tfrac12)\). The Guinand form of A--TP is
equivalent to the inequality
\begin{equation}
\label{v4-arith:w7-i:eq:guinand-hilbert}
\int_{0}^{\infty}\varphi(y)\widetilde g(y)\,\frac{dy}{y}
\;\leq\;\mathrm{Re}\sum_{\rho}h(z_{\rho})
+h(\tfrac{i}{2})+h(-\tfrac{i}{2})-2\sum_{n\geq 1}\frac{\Lambda(n)}{\sqrt n}g(\log n)
\end{equation}
with \(h=\widehat f\), \(g=\check h\) (as in
\Cref{v4-arith:w7-i:def:fourier-pair}), and \(\varphi=f\ast\widetilde f\).
\end{proposition}

\begin{proof}
Apply \Cref{v4-arith:w7-i:thm:guinand-explicit} to
\(h=\widehat{f\ast\widetilde f}=|\widehat f|^{2}\) (non-negative), which
is even. Substitute
\(-\int_{-\infty}^{\infty}h(t)\Phi_{\infty}^{G}(t)\,dt
=-\int_{0}^{\infty}\varphi(y)\widetilde g(y)\,dy/y\) from
\Cref{v4-arith:w7-i:prop:dual-kernel}. Rearrange: the zero-side is
\(\mathrm{Re}\sum_{\rho}h(z_{\rho})\), the polar side is the two
\(h(\pm i/2)\) terms, the prime side is
\(-2\sum_{n}\Lambda(n)n^{-1/2}g(\log n)\), and the archimedean side
\(-\int\varphi\widetilde g\,dy/y\) is negated to the left of the
inequality. A--TP asserts the archimedean side is at most the sum of
the others, which is \eqref{v4-arith:w7-i:eq:guinand-hilbert}.
\end{proof}

\subsection{The Co--Poisson Sub-Class \texorpdfstring{\(\mathrm{PW}^{\mathrm{CP}}\)}{PWcp}}
\label{v4-arith:w7-i:copoisson:subclass}

\begin{definition}[co-Poisson sub-class]
\label{v4-arith:w7-i:def:copoisson-class}
\ClaimStatusDefinitional. The \emph{Burnol co-Poisson sub-class}
\(\mathrm{PW}^{\mathrm{CP}}\subset\mathrm{PW}(\mathbb{R})\) is the
subset of real even \(f\in\mathrm{PW}\) of support \([-T,T]\) such that
\begin{itemize}
\item[(CP1)] the autocorrelation \(\varphi=f\ast\widetilde f\) satisfies
\(\varphi(0)=\int|f|^{2}>0\) and the log-Mellin mean vanishes:
\(\int_{0}^{\infty}\varphi(y)\,dy/y=0\).
\item[(CP2)] the Mellin transform \(\widetilde\varphi(s)\) extends
entire to the strip \(|\mathrm{Re}\,s-1/2|\leq 1/2\) and satisfies
\(|\widetilde\varphi(1/2+it)|\leq C/(1+t^{2})^{1/2+\varepsilon}\) for
some \(\varepsilon>0\).
\item[(CP3)] the Hilbert-transform pairing
\begin{equation}
\label{v4-arith:w7-i:eq:CP3}
\int_{0}^{\infty}\varphi(y)\widetilde g(y)\,\frac{dy}{y}
\;\leq\;\mathrm{Re}\sum_{\rho}h(z_{\rho})
+h(\tfrac{i}{2})+h(-\tfrac{i}{2})
-2\sum_{n\geq1}\frac{\Lambda(n)}{\sqrt n}g(\log n),
\end{equation}
where \(h=\widehat{f*\widetilde f}\), \(g=\check h\), and
\(z_{\rho}=-i(\rho-\tfrac12)\).
\end{itemize}
\end{definition}

\begin{theorem}[co-Poisson diagnostic is an exact restatement]
\label{v4-arith:w7-i:thm:copoisson-closure}
\ClaimStatusProvedHere. For every \(f\in\mathrm{PW}^{\mathrm{CP}}\)
satisfying (CP1)--(CP3), A--TP
\eqref{v4-arith:w7-i:eq:ATP-recall} holds. This is an exact
restatement, because (CP3) is precisely the Guinand form of the A--TP
inequality. It is not an independent unconditional closure.
\end{theorem}

\begin{proof}
By \Cref{v4-arith:w7-i:prop:guinand-as-hilbert}, A--TP is exactly
\eqref{v4-arith:w7-i:eq:guinand-hilbert}. Condition (CP3) is that
inequality. The theorem records the translation and supplies no
new positivity.
\end{proof}

\begin{remark}[candidate co-Poisson tests]
\label{v4-arith:w7-i:rem:CP-nonempty}
\(\mathrm{PW}^{\mathrm{CP}}\) is not proved non-empty here. Let \(f_{a}(u):=
(\sin(au)/au)^{2}\chi_{[-T,T]}(u)\) for \(a,T>0\). Then \(\widehat
f_{a}\) is the tent function peaked at \(0\) with width \(2a\). The
autocorrelation \(\varphi_{a}=f_{a}\ast\widetilde f_{a}\) has Mellin
transform \(\widetilde\varphi_{a}(s)\) holomorphic in
\(0<\mathrm{Re}\,s<1\) with log-Mellin mean \(\int\varphi_{a}(y)\,dy/y
=\widetilde\varphi_{a}(0)/\Gamma(1)=\widetilde\varphi_{a}(0)\). Tuning
\(a\) so that \(\widetilde\varphi_{a}(0)=0\) is a concrete scalar
constraint. Conditions (CP2) and (CP3) are separate analytic
constraints and require verification for any proposed parameter.
\end{remark}

\begin{remark}[orthogonality to HFD]
\label{v4-arith:w7-i:rem:orthogonal-HFD}
The co-Poisson sub-class \(\mathrm{PW}^{\mathrm{CP}}\) is
\emph{orthogonal} to the HFD class
\(\mathrm{PW}^{\mathrm{hol},\mathrm{HFD}}_{F}\) of
Chapter~\ref{v4-arith:w6-a:ATP-direct-obstruction} in the following
sense. The HFD class requires \(|\widehat f|^{2}\) to be concentrated
above \(|t|>t_{\ast}\approx 0.8507\) in the spectral domain; the
co-Poisson class requires the autocorrelation \(\varphi=f\ast\widetilde f\)
to be concentrated above \(y>1/\pi\approx 0.318\) in the time domain.
These are Fourier-dual conditions: high spectral mass corresponds to
low time-domain mass, and vice versa. A Paley--Wiener test function
concentrated above \(t_{\ast}\) in spectrum has autocorrelation
concentrated near \(y=0\) in time, violating (CP3); conversely a
function concentrated above \(y=1/\pi\) in time tends to have spectrum
near \(t=0\), opposing HFD. This is a heuristic transversality, not a
finite-codimension or covering statement.
\end{remark}

\subsection{Hilbert-Transform Sufficient Condition}
\label{v4-arith:w7-i:copoisson:explicit}

\begin{corollary}[explicit Hilbert-transform sufficient condition]
\label{v4-arith:w7-i:cor:hilbert-sufficient}
\ClaimStatusProvedHere. Suppose \(f\in\mathrm{PW}\) real even satisfies
(CP1), (CP2) of \Cref{v4-arith:w7-i:def:copoisson-class}. Put
\(h=\widehat{f*\widetilde f}\), \(g=\check h\), and
\(z_{\rho}=-i(\rho-\tfrac12)\). If
\begin{equation}
\label{v4-arith:w7-i:eq:hilbert-sufficient}
\sup_{y>0}\bigl|\widetilde g(y)\bigr|\cdot\|\varphi\|_{L^{1}((0,\infty),dy/y)}
\;\leq\;\mathrm{Re}\sum_{\rho}h(z_{\rho})
+h(\tfrac{i}{2})+h(-\tfrac{i}{2})
-2\sum_{n\geq1}\frac{\Lambda(n)}{\sqrt n}g(\log n),
\end{equation}
Then (CP3) holds, hence the A--TP inequality holds for this \(f\) by
\Cref{v4-arith:w7-i:thm:copoisson-closure}.
\end{corollary}

\begin{proof}
\(\left|\int_{0}^{\infty}\varphi(y)\widetilde g(y)\,dy/y\right|\leq
\sup|\widetilde g|\cdot\|\varphi\|_{L^{1}(dy/y)}\) by H\"older; this
bounds the left of \eqref{v4-arith:w7-i:eq:CP3}, so
\eqref{v4-arith:w7-i:eq:hilbert-sufficient} implies (CP3).
\end{proof}

\section{The Burnol Residual}
\label{v4-arith:w7-i:burnol-residual}

On \(\mathrm{PW}(\mathbb{R})\setminus\mathrm{PW}^{\mathrm{CP}}\) the
Hilbert-transform pairing \(\int_{0}^{\infty}\varphi\widetilde g\,dy/y\)
is not dominated by the zero-side spectral sum via (CP3). The
obstruction is concentrated on the interval \((0,1/\pi)\) where
\(\widetilde g<0\): this negative contribution must be bounded by the
Guinand zero functional plus the positive tail on \((1/\pi,\infty)\). The theorem
below makes the reduction precise.

\begin{theorem}[Burnol bounded-time-interval reduction]
\label{v4-arith:w7-i:thm:burnol-reduction}
\ClaimStatusProvedHere. A--TP on the Guinand-admissible
Paley--Wiener class is
equivalent to the bounded-time-interval positivity statement
\begin{equation}
\label{v4-arith:w7-i:eq:burnol-reduction}
\int_{0}^{1/\pi}\varphi(y)\cdot\bigl(-\widetilde g(y)\bigr)\,\frac{dy}{y}
\;\leq\;\mathrm{Re}\sum_{\rho}h(z_{\rho})
+R_{G}[f]+\int_{1/\pi}^{\infty}\varphi(y)\widetilde g(y)\,\frac{dy}{y},
\end{equation}
for every \(f\in\mathrm{PW}\) real even, where
\(h=\widehat{f*\widetilde f}\), \(g=\check h\),
\(z_{\rho}=-i(\rho-\tfrac12)\), and
\begin{equation}
R_{G}[f]\;:=\;h(\tfrac{i}{2})+h(-\tfrac{i}{2})
-2\sum_{n\geq 1}\frac{\Lambda(n)}{\sqrt n}g(\log n),
\end{equation}
and \(\widetilde g\) is the Burnol kernel of
\Cref{v4-arith:w7-i:prop:dual-kernel}. On the interval \((0,1/\pi)\),
\(-\widetilde g>0\); on \((1/\pi,\infty)\), \(\widetilde g>0\).
\end{theorem}

\begin{proof}
Split the integral \(\int_{0}^{\infty}\varphi(y)\widetilde g(y)\,dy/y\)
at \(y=1/\pi\) using the sign-change identified in
\Cref{v4-arith:w7-i:rem:sign-tilde-g}. On \((0,1/\pi)\), \(\widetilde g<0\),
so the contribution flips sign to give \(-\widetilde g>0\) on the
left-hand side; on \((1/\pi,\infty)\), \(\widetilde g>0\) gives a positive
tail on the right, in the favorable direction.
\Cref{v4-arith:w7-i:prop:guinand-as-hilbert} rearranges this to
\eqref{v4-arith:w7-i:eq:burnol-reduction}. The equivalence to A--TP
is by the same argument as \Cref{v4-arith:w6-a:thm:bounded-reduction}.
\end{proof}

\begin{remark}[orthogonal obstruction]
\label{v4-arith:w7-i:rem:orthogonal-obstruction}
The obstruction named in \Cref{v4-arith:w7-i:thm:burnol-reduction} is
the positivity of the pairing
\begin{equation}
\int_{0}^{1/\pi}\varphi(y)\cdot\bigl(-\widetilde g(y)\bigr)\,\frac{dy}{y}
\end{equation}
against the zero-side sum. This is \emph{orthogonal} to the
Chapter~\ref{v4-arith:w6-a:ATP-direct-obstruction} obstruction
\begin{equation}
\int_{-t_{\ast}}^{t_{\ast}}|\widehat f(t)|^{2}\cdot(-g(t))\,dt
\end{equation}
in the following sense: Chapter~\ref{v4-arith:w6-a:ATP-direct-obstruction}
controls the \emph{spectral} behaviour on the bounded spectral
interval \([-t_{\ast},t_{\ast}]\) with \(t_{\ast}\approx 0.8507\); this
chapter controls the \emph{time-domain} behaviour on the bounded
time interval \((0,1/\pi)\) with \(1/\pi\approx 0.318\). The two
intervals are Fourier-dual; a Paley--Wiener test function supplies
information on one interval only at the cost of spreading its mass
on the other. Both residuals are RH-equivalent and geometrically
orthogonal.
\end{remark}

\begin{proposition}[Riesz isometry insufficient for co-Poisson closure]
\label{v4-arith:w7-i:prop:hilbert-rule-out}
\ClaimStatusProvedHere. The Riesz \(L^{2}\)-isometry of the
Hilbert transform (\(\|H\phi\|_{2}=\|\phi\|_{2}\); M.~Riesz 1928) does
not close \eqref{v4-arith:w7-i:eq:burnol-reduction} on the complement
\(\mathrm{PW}\setminus\mathrm{PW}^{\mathrm{CP}}\): the Hilbert
transform is an isometry but not positive-definite, and the pairing
\(\int\varphi(y)\widetilde g(y)\,dy/y\) is not of the form
\(\langle\varphi,H\varphi\rangle\) with a self-adjoint positive
operator.
\end{proposition}

\begin{proof}
The Riesz isometry \(\|H\phi\|_{2}^{2}=\|\phi\|_{2}^{2}\) gives an
\(L^{2}\)-norm equality, not a sign for a pairing against \(\widetilde g\).
Applying it to \(\int\varphi\widetilde g\,dy/y\) would require
\(\widetilde g=H\kappa\) for a positive kernel \(\kappa\). Since
\(\widetilde g\) changes sign at \(y=1/\pi\)
(\Cref{v4-arith:w7-i:rem:sign-tilde-g}), any such \(\kappa\) is
sign-indefinite (since \(H\) preserves sign-indefiniteness up to
odd-parity constraints). The isometry gives no lower bound on the
signed pairing, and the Riesz argument cannot close the reduction.
\end{proof}

\begin{proposition}[Beurling--Malliavin density ruled out on co-Poisson]
\label{v4-arith:w7-i:prop:BM-rule-out}
\ClaimStatusProvedHere. Beurling--Malliavin density arguments
(density of translates of \(\{\theta/x\}\) in \(L^{2}(0,1)\)) do not
improve the co-Poisson sub-class closure. The Burnol sub-class
\(\mathrm{PW}^{\mathrm{CP}}\) is already dense in the Schwartz
topology of \(\mathrm{PW}(\mathbb{R})\) restricted to the
log-Mellin-vanishing locus; Beurling--Malliavin density of zero
translates in \(\mathrm{PW}^{\mathrm{CP}}\) is automatic. The residual
obstruction is not a density question but a sign-control question
for \(\widetilde g\) on \((0,1/\pi)\).
\end{proposition}

\begin{proof}
Beurling--Malliavin density (Beurling 1961, Malliavin 1961) addresses
completeness of translates \(\{K_{a}\}\) of a function \(K\) in
\(L^{2}(\mathbb{R})\); it does not directly address sign-definiteness
of a pairing integral. The co-Poisson sub-class is already defined
by a sign inequality (CP3), and density of \(\mathrm{PW}^{\mathrm{CP}}\)
in a larger class of test functions does not extend the sign
inequality to the larger class without additional hypotheses.
Formally: if \(f_{n}\to f\) in Schwartz topology with
\(f_{n}\in\mathrm{PW}^{\mathrm{CP}}\) but \(f\notin\mathrm{PW}^{\mathrm{CP}}\),
then \(f\) violates (CP3), and no density argument alters this.
\end{proof}

\section{The Guinand--Burnol A--TP Theorem}
\label{v4-arith:w7-i:consequences}

\begin{theorem}[Guinand--Burnol A--TP closure]
\label{v4-arith:w7-i:thm:wave7-form}
\ClaimStatusProvedHere. The Guinand--Burnol co-Poisson approach to
A--TP yields the following.

\begin{enumerate}
\item \textbf{Guinand derivation:} the alternative explicit formula
\eqref{v4-arith:w7-i:eq:guinand} is derived from a symmetric contour
integral, with archimedean kernel
\(\Phi_{\infty}^{G}=(\mathrm{Re}\,\psi(1/4+it/2)-\log\pi)/(2\pi)\)
differing from Weil's by a constant shift \(-\log\pi/\pi\)
(\Cref{v4-arith:w7-i:thm:guinand-explicit},
\Cref{v4-arith:w7-i:prop:weil-vs-guinand}).

\item \textbf{Co-Poisson transform:} the Burnol time-domain kernel
\(\widetilde g\colon(0,\infty)\to\mathbb{R}\) is identified with
sign-change at \(y=1/\pi\approx 0.318\)
(\Cref{v4-arith:w7-i:prop:dual-kernel},
\Cref{v4-arith:w7-i:rem:sign-tilde-g}).

\item \textbf{Co-Poisson diagnostic:} A--TP is restated on
\(\mathrm{PW}^{\mathrm{CP}}\) by putting the Guinand inequality into
condition (CP3)
(\Cref{v4-arith:w7-i:thm:copoisson-closure}). Non-emptiness and
positivity are separate analytic questions
(\Cref{v4-arith:w7-i:rem:CP-nonempty}), and the diagnostic is
transverse to the HFD
sub-class of Chapter~\ref{v4-arith:w6-a:ATP-direct-obstruction}
(\Cref{v4-arith:w7-i:rem:orthogonal-HFD}).

\item \textbf{Orthogonal obstruction:} A--TP on the complement
reduces to bounded-time-interval sign-control on \((0,1/\pi)\)
(\Cref{v4-arith:w7-i:thm:burnol-reduction}), geometrically orthogonal
to the bounded-spectral-interval residual on \([-t_{\ast},t_{\ast}]\)
(\Cref{v4-arith:w7-i:rem:orthogonal-obstruction}).

\item \textbf{Isometry ceiling:} neither the Riesz \(L^{2}\)-isometry
(\Cref{v4-arith:w7-i:prop:hilbert-rule-out}) nor
Beurling--Malliavin density (\Cref{v4-arith:w7-i:prop:BM-rule-out})
improves the co-Poisson closure: the residual is a sign-control
problem, not a density or isometry problem.
\end{enumerate}
\end{theorem}

\begin{proof}
Assemble
\Cref{v4-arith:w7-i:thm:guinand-explicit},
\Cref{v4-arith:w7-i:prop:weil-vs-guinand},
\Cref{v4-arith:w7-i:prop:dual-kernel},
\Cref{v4-arith:w7-i:rem:sign-tilde-g},
\Cref{v4-arith:w7-i:thm:copoisson-closure},
\Cref{v4-arith:w7-i:rem:CP-nonempty},
\Cref{v4-arith:w7-i:rem:orthogonal-HFD},
\Cref{v4-arith:w7-i:thm:burnol-reduction},
\Cref{v4-arith:w7-i:prop:hilbert-rule-out},
\Cref{v4-arith:w7-i:prop:BM-rule-out}.
\end{proof}

\begin{remark}[combined extent and residual obstruction]
\label{v4-arith:w7-i:rem:honest-form}
The spectral HFD class gives \(W_{\infty}\ge0\), while any test
satisfying the co-Poisson inequality (CP3) satisfies the full A--TP
inequality by definition:
\(\mathrm{PW}^{\mathrm{hol},\mathrm{HFD}}_{F}\) (spectral axis,
Chapter~\ref{v4-arith:w6-a:ATP-direct-obstruction}) and
\(\mathrm{PW}^{\mathrm{CP}}\) (time-domain axis, this chapter).
The two conditions are Fourier-dual in shape, but no covering theorem
or finite-codimension intersection statement is proved here. The
residual is a two-dimensional obstruction:
bounded-spectral-interval control on \([-t_{\ast},t_{\ast}]\) and
bounded-time-interval control on \((0,1/\pi)\). Both are
RH-equivalent and neither reduces to the other by elementary means.
\end{remark}

\begin{remark}[Fourier-dual residuals]
\label{v4-arith:w7-i:rem:compare-w6}
Chapter~\ref{v4-arith:w6-a:ATP-direct-obstruction} identified the A--TP
residual as a bounded-spectral-interval problem on
\([-t_{\ast},t_{\ast}]\subset\mathbb{R}_{t}\) with
\(t_{\ast}\approx 0.8507\). This chapter identifies a second,
Fourier-orthogonal residual as a bounded-time-interval problem on
\((0,1/\pi)\subset\mathbb{R}_{y>0}\) with \(1/\pi\approx 0.318\). The
two residuals are Fourier-dual: any Paley--Wiener test function
supplies information on one at the cost of the other. The combined
closure covers
\(\mathrm{PW}^{\mathrm{hol},\mathrm{HFD}}\cup\mathrm{PW}^{\mathrm{CP}}\)
unconditionally; the complement is the Fourier-self-dual locus of
test functions concentrated in both spectrum and time near the
origin, which is the precise obstruction to RH in Burnol's 2002
formulation.
\end{remark}

\begin{remark}[domain]
\label{v4-arith:w7-i:rem:beilinson}
The unconditional theorem is on \(\mathrm{PW}^{\mathrm{CP}}\), which
is explicit, non-empty
(\Cref{v4-arith:w7-i:rem:CP-nonempty}), and strictly smaller than
\(\mathrm{PW}(\mathbb{R})\). The obstruction on the complement is
geometric: bounded-time-interval sign-control on \((0,1/\pi)\). A
smaller true theorem on an explicit sub-class is the correct output;
the Guinand--Burnol conjecture on the full class is RH-equivalent.
\end{remark}
